\section{Specification-Traceable Verification}
\label{sec:verification}

The verification structure in Fig.~\ref{fig:verification} separates stimulus,
observation, prediction, assertions, and coverage accounting. The separation is
important for training logic: a scoreboard that predicts from the DUT's
internal LFSR state or phase register can agree with the same defect it is meant
to detect.

\begin{figure*}[t]
\centering
\includegraphics[width=0.86\textwidth]{fig07_verification.pdf}
\caption{Verification architecture. The predictor reconstructs message order,
per-lane LFSR state, thresholds, and error cause from interface observations.}
\label{fig:verification}
\end{figure*}

\subsection{Directed and Randomized Layers}

Directed tests isolate handshake and boundary semantics before integration.
The sideband suite covers successful exchange, timeout/retry, retry exhaustion,
response mismatch, and abort. LFSR tests cover nonzero seeds, the implemented
polynomial, receive gaps, strict threshold equality, abort, the default
4096-sample configuration, and 16-bit saturation. Error-manager tests cover
cause priority, retention, immediate/acknowledged/timed entry, protected clear,
and count saturation. Wrapper tests read version, state, status, training phase,
cause/count, allowed and ignored clear operations, and invalid addresses.

Nine deterministic UVM tests retain top-level state and transition checks.
Separate seeded campaigns stress sideband, training, and recovery behavior.
The production integrated test disables the abstract CENTER1 bypass and uses
nine scenario classes: nominal completion; pattern failure followed by retry
and pass; START retry; malformed START response; END retry; malformed END
response; and reset, fatal, or state-timeout aborts. Backpressure, response
latency, sample gaps, corruption, threshold, and abort phase vary within these
classes.

\subsection{Independent Predictor}

The predictor observes transmitted requests, received responses, accepted
training samples, and public controller/error status. It maintains its own
23-bit state for each of 16 lanes, uses the specified project seed table and
polynomial, advances only when a receive sample is accepted, and computes the
expected mismatch count independently. It also maintains a protocol phase:
START must close before pattern activity, END must not appear before a passing
pattern, and DATATRAINVREF must not be reached until END closes. Abort scenarios
check both clean cessation and retained cause attribution.

\subsection{Assertions and Coverage Boundary}

Concurrent assertions check that \texttt{link\_up} is exclusive to ACTIVE,
timeout routes to TRAINERROR, reset cannot jump directly to ACTIVE, and states
remain known. Integrated assertions restrict pattern activity to
DATATRAINCENTER1/PATTERN, constrain requests to their phase, require an END
response before production-mode advance, check the strict error threshold, and
hold LFSR progress when no sample is accepted. Recovery properties cover
pending request retention, bounded TRAINERROR entry, pending clear, log
stability, and ignored clear in TRAINERROR.

Explicit counters record scenario, phase, gap, corruption, threshold, retry,
abort, and selected cross combinations. They demonstrate that planned classes
were sampled, but are not reported as a standard functional-coverage percentage
because no UCDB was produced under the available Starter license.
